\documentclass[a4paper,11pt,dvipdfmx]{jsarticle}
\usepackage{amsmath, amssymb, amsthm}
\usepackage{bm}
\usepackage{mathtools}
\usepackage{comment}
\usepackage{algorithm}
\usepackage{algorithmic}
\usepackage{bm}
\usepackage{mathtools}
\usepackage{enumitem}
\mathtoolsset{showonlyrefs}
\usepackage[dvipdfmx, colorlinks=true, linkcolor=blue, citecolor=blue, urlcolor=blue, setpagesize=false]{hyperref}
\usepackage{pxjahyper}
% --- 定理・定義環境 ---
\theoremstyle{definition}
\newtheorem{definition}{定義}[section]
\newtheorem{theorem}[definition]{定理}
\newtheorem{lemma}[definition]{補題}
\newtheorem{assumption}{仮定}
\newtheorem{proposition}[definition]{命題}
\newtheorem{requirement}[definition]{要請}
\newtheorem{example}[definition]{例}
\newcommand{\fU}{\underline{f}}
\newcommand{\gU}{\underline{g}}
\newcommand{\uU}{\underline{u}}
\newcommand{\aU}{\underline{a}}
\newcommand{\bU}{\underline{b}}
\newcommand{\cU}{\underline{c}}
\newcommand{\dU}{\underline{d}}
\newcommand{\vU}{\underline{v}}
\newcommand{\zeroU}{\underline{0}}
\newcommand{\C}[2]{C_{{#1},{#2}}}
\newcommand{\cy}[3]{\mathcal{C}_{{#1},{#2},{#3}}}
\newcommand{\inv}{^{-1}}
\newcommand{\HHX}{\hat{H}_X}
\newcommand{\HHZ}{\hat{H}_Z}
\newcommand{\tr}{^\top}

% =========================================
\begin{document}

\title{}
\author{}
\date{\today}
\maketitle
\tableofcontents
\clearpage
\section{02/04のゼミの内容}
\begin{itemize}
  \item 長さ6のサイクルがあるのではないか：関数$f(x)$と行列$F$の関係について、\cite{kasai2025quantum}とは行と列が逆、アクティブ部でのみサイクルを考える。
  \item UTCBC以外の長さ8のサイクルを避けるべきか？：今は長さ6以外のサイクルを含まないことを条件にしているが、もしUTCBC以外のサイクルも避ければ、性能が良くなるかも？(ならないかも)、エラーがどのサイクル由来かを調べられるとよい
  \item 整数計画法で求めた特殊解はランダムネスが小さい
  \item APMを使っている時点でランダムネスは減少しているのでAPMに縛らないとランダムネスを増やせる
  \item 他の改良案：Pを小さく(2,3以外の素因数を使う)、ガースを大きく(アクティブ行列の取り出し方を変える。)
\end{itemize}

\section{関数と行列の関係}
\cite{kasai2025quantum}では、
\[
f(j) = i \iff F_{i,j} = 1
\]
\cite{kasai2026breakingorthogonalitybarrierquantum}では、
\[
f(i) = j \iff F_{i,j} = 1
\]
よって、サイクルの合成関数も
サイクル$c=f_0 \rightarrow f_1 \rightarrow\dots\rightarrow f_n$ならば
\[
f_c(x) = f_n f_{n-1}\inv\dots f_1 f_{0}\inv(x)
\]

\section{UTCBCのカウント}
UTCBCを構成する$\rU = [r_0, r_1, r_2, r_3]\in[J]^4, \cU = [c_0, c_1, c_2, c_3]\in[L]^4$は以下の条件を満たす。
\begin{align}
  &c_2 - r_1 \equiv c_0 - r_0\pmod{L/2}\label{02061}\\
  &c_3 - r_2 \equiv c_1 - r_1\pmod{L/2}\label{02062}\\
  &c_3 - r_3 \equiv c_1 - r_0\pmod{L/2}\label{02063}\\
  &r_0\neq r_1, r_1\neq r_2, r_2\neq r_3, r_3\neq r_0\label{02064}\\
  &c_0\neq c_1, c_1\neq c_2, c_2\neq c_3, c_3\neq c_0\label{02065}\\
  &(c_0\in[L/2]\land c_2\in[L/2])\lor (c_0\in[L/2, L]\land c_2\in[L/2, L])\label{02071}\\
  &(c_1\in[L/2]\land c_3\in[L/2])\lor (c_1\in[L/2, L]\land c_3\in[L/2, L])\label{02072}\\
\end{align}
式\eqref{02062}, \eqref{02063}を整理すると
\begin{align}
  &c_3 - c_1 \equiv r_2 - r_1\pmod{L/2}\\
  &c_3 - c_1 \equiv r_3 - r_0\pmod{L/2}\\
\end{align}
なので、
\begin{align}
  &c_2 - c_0 \equiv r_1 - r_0(=\delta_1\neq 0)\pmod{L/2}\\\
  &c_3 - c_1 \equiv r_2 - r_1 \equiv r_3 - r_0(=\delta_2\neq 0)\pmod{L/2}\label{02066}
\end{align}
が成り立つ。
これを満たす$\rU$は、
\begin{align}
  [0, 1, 2, 1]&((\delta_1, \delta_2) = (1, 1))\\
  [1, 0, 1, 2]&((\delta_1, \delta_2) = (-1, 1))\\
  [1, 2, 1, 0]&((\delta_1, \delta_2) = (1, -1))\\
  [2, 1, 0, 1]&((\delta_1, \delta_2) = (-1, -1))\\
\end{align}
のいずれかである。

これらはすべて$[0, 1, 2, 1]$の左シフトで表せるので、$\rU = [0,1,2,1]$とする。
$\cU$の組み合わせを考える。
$(\delta_1, \delta_2) = (1, 1)$なので、
\begin{align}
  &c_2 - c_0 \equiv \delta_1\pmod{L/2}\\
  &c_3 - c_1 \equiv \delta_2\pmod{L/2}
\end{align}
であり、$c_0$を決めると、$c_2=c_0+\delta_1$が決まり、$c_0, c_2$と異なるように$c_1$を決めると、$c_3=c_1+\delta_2$が決まる。

ここで、式\eqref{02065}を考慮する。
$c_3$が$c_0$または$c_2$と重複してしまう$c_1$が1つ存在する。
さらに、逆順に辿ることものが存在することを考えると、$\cU$の組み合わせは
$12\times (12-3)/2 = 54$通り。

\section{検査するサイクルの再考}
段階的に目標を達成するために、まずは長さ6以下のサイクルを含まない$\fU, \gU$の一般解の導出を目指す。
\paragraph{長さ4のサイクル}
長さ4のサイクルは、2つの位置を与えれば、構成する4つの位置が決まる。
2つの位置の選び方は、列、行ともに異なるものを選ぶ必要があることを考えると、1つ目が$J\times L$通り、2つ目は$(J-1)\times (L-1)$通りであり、全部で$L/2\times (L/2-1)\times L\times (L-1)$通りである。
よって、検査すべきサイクルの個数は
\begin{align}
  \frac{J\times (J-1)\times L\times (L-1)}{4}\label{eq:012803}
\end{align}
個である。
$J=3, L=12$を代入すると、
\begin{align}
  \frac{3\times 2\times 12\times 11}{4} = 198
\end{align}
個である。

\paragraph{長さ6のサイクル}
長さ6のサイクルは、3つの位置を与えれば、サイクルが一意に決まる。
3つの位置の選び方は、行インデックス$[r_0,r_1,r_2]$と列インデックス$[c_0,c_1,c_2]$がそれぞれ相異なる必要があるので、
\begin{align}
  J\times (J-1)\times (J-2)\times L\times (L-1)\times (L-2)
\end{align}
通りである。
また、同一のサイクルを与える3つの位置の並べ方は、サイクルの開始点の選び方(3通り)と巡回方向(2通り)により$6$通りある。
したがって、検査すべきサイクルの個数は
\begin{align}
  \frac{J\times (J-1)\times (J-2)\times L\times (L-1)\times (L-2)}{6}\label{eq:012805}
\end{align}
個である。
$J=3, L=12$を代入すると、
\begin{align}
  \frac{3\times 2\times 1\times 12\times 11\times 10}{6} = 1320
\end{align}
個である。

以上合わせると、検査すべきサイクルは
$198+1320 = 1518$個であり、$\HHX, \HZ$で考えるので、制約の数は3136個となる。

\section{ガースが6になるような$F_i, G_i$の一般解を求める}
まず、$F_i, G_i$が満たすべき条件は以下である。
\begin{enumerate}[label=\Alph*, ref=\Alph*]
  \item $f_i,G_j$は$(i,j)\in[L/2]^2\backslash\{(0,3), (1,2)\}$に対して可換である\label{cond:a}
  \item 少なくとも一つの$(i,j)\in\{(0,3), (1,2)\}$に対して非可換である\label{cond:b}
  \item アクティブ行列$\HX,\HZ$内の短いサイクルを避ける(長さ6以下のサイクル)\label{cond:c}
\end{enumerate}

\begin{definition}\label{def:02081}
  行列$G$を以下の$L,R$を用いて$G = [L\mid R]$と定義する。
\begin{align}
  L_{(6i+j, k)} &= \begin{cases}
    1-a_{g_i} & \mif k = i\\
    0 & \text{otherwise}
  \end{cases}\\
  R_{(6i+j, k)} &= \begin{cases}
    a_{f_j}-1 & \mif k = j\\
    0 & \text{otherwise}
  \end{cases}
\end{align}
\end{definition}

\begin{definition}\label{def:02082}
  $G$から3,8行目($(i,j)=(0,3),(1,2)$に対応する行)を除いたものを$G_a$とし、3,8行目のみを取り出したものを$G_b$とする。
\end{definition}

\subsection{条件\ref{cond:a}}
条件\ref{cond:a}は、
\begin{align}
  G_a \bU \equiv \zeroU\label{02083}
\end{align}
で表せ、これをスミス標準形を用いて解くと、解空間
\begin{align}
  \bU = \sum_{i=0}^{d_a} c_i\vU_i\label{02084}
\end{align}
が得られる($d_a$は解空間の次元)。
ここで、$V=[\vU_0\tr,\dots,\vU_{d_a}\tr]\quad \cU=[c_0,\dots,c_{d_a}]\tr$とすると、
\begin{align}
  \bU = V\cU\label{02085}
\end{align}
と表せる。\\
$\rightarrow$この記述ですべての解空間を網羅できるわけではないことが分かった。
スミス標準形を用いて解空間を求めると、条件を満たさないものも含む空間の基底が出てくる。
一方でエルミート標準形を用いて解空間を求めると、条件を満たすものだけが含まれる空間の基底が出てくるが、この空間に含まれない解も存在する。(論文で挙げられた解は含まれていない)

\subsection{条件\ref{cond:b}}
条件\ref{cond:b}は、
\begin{align}
  G_b \bU \equiv \dU_b\label{02086}
\end{align}
かつ$\dU_b$のすべての要素が非ゼロ、で表せる。
式\eqref{02085}を代入すると、
\begin{align}
  &\qquad G_b V\cU \equiv \dU_b\label{02087}\\
  &\iff G_b'\cU \equiv \dU_b\label{02088}\\
  &\iff {G_b'}_{(k)}\cU \not\equiv 0\text{ for all }k \in \{0,1\}\label{02089}
\end{align}
が得られる。ただし、$G_b' = G_b V$であり、$G_{(i)}$は行列$G$の$i$行目である。
これにより、各$k$に対して$\cU$の禁止領域が求められる。

\subsection{条件\ref{cond:c}}
条件\ref{cond:c}は、34530個ものサイクルについて考える必要がある。合成関数が$f_C(x) = a_C x + b_C$であるサイクルを考える。
条件は、
\begin{align}
  b_C \not\equiv 0 \pmod{\gcd(a_C-1, P)}\label{021101}
\end{align}
が成り立つことで、$b_C$は$\bU$の線形結合として$b_C = c_C \bU$で表せる。
これを用いると、\eqref{021101}は以下のように表せる。
\begin{align}
  c_C\bU \not\equiv 0 \pmod{\gcd(a_C-1, P)}\label{021102}\\
  \iff \frac{P}{\gcd(a_C-1, P)}c_C\bU \not\equiv 0 \pmod{P}\label{021103}
\end{align}
$\bU = V\cU$より、
\begin{align}
  c_C\bU \not\equiv 0 \pmod{\gcd(a_C-1, P)}\label{021104}\\
  \iff \frac{P}{\gcd(a_C-1, P)}c_C\bU \not\equiv 0 \pmod{P}\label{021105}\\
  \iff \frac{P}{\gcd(a_C-1, P)}c_C V\cU \not\equiv 0 \pmod{P}\label{021106}\\
  \iff c_C'\cU \not\equiv 0 \pmod{P}\label{021107}
\end{align}
ただし、$c_C' = \frac{P}{\gcd(a_C-1, P)}c_C V$である。

条件\ref{cond:b},\ref{cond:c}合わせて、69062個の制約が得られる。

\section{係数ベクトル $\underline{b}$ の一般解の定式化と探索方針}

本章では、量子LDPC符号の性能を左右するアフィン置換行列の平移係数ベクトル $\underline{b}$ について、直交性制約とガース（girth）制約を同時に満たす一般解を代数的に特定する手法を詳述する。

\section{中国剰余定理を用いた解空間の分解と一般解の代数的構成}

本節では、探索空間の巨大化に伴う計算コストの問題を解決するため、中国剰余定理（Chinese Remainder Theorem, CRT）を用いた解空間の分解手法を提案する。本手法は、単なる全探索（しらみつぶし）とは異なり、代数的構造を利用して効率的に一般解を特定するものである。

\subsection{法 $P$ の分解と成分表示}
本研究における法 $P=768$ は、$P = 3 \times 2^8 = 3 \times 256$ と素因数分解される。
CRTにより、環 $\mathbb{Z}_{768}$ は以下の直積環に同型である。
\begin{equation}
    \mathbb{Z}_{768} \cong \mathbb{Z}_3 \times \mathbb{Z}_{256}
\end{equation}
これに伴い、係数ベクトル $\mathbf{x} \in \mathbb{Z}_{768}^{d_a}$ は、各成分の組 $(\mathbf{x}^{(3)}, \mathbf{x}^{(256)})$ によって一意に決定される。

\subsection{禁止領域回避の十分条件}
禁止ベクトル $\mathbf{r}_i$ に対する非零条件 $\mathbf{r}_i^T \mathbf{x} \not\equiv 0 \pmod{768}$ を満たすためには、以下の条件が少なくとも一領域で成立すれば十分である。
\begin{equation}
    \mathbf{r}_i^T \mathbf{x} \not\equiv 0 \pmod 3 \quad \text{or} \quad \mathbf{r}_i^T \mathbf{x} \not\equiv 0 \pmod{256}
\end{equation}
特に $\mathbb{Z}_3$ は体であるため、超平面配置の回避問題としての性質が顕著であり、この成分を適切に選択することで、法 $768$ における禁止領域の大部分を代数的に排除できる。

\subsection{構成的探索アルゴリズム}
既知の特殊解を $\mathbf{x}_0$ とし、その成分を $\mathbf{x}_0 \equiv (\mathbf{x}_0^{(3)}, \mathbf{x}_0^{(256)})$ とする。一般解の集合を効率的に得るため、以下の手順で探索を行う。

\begin{enumerate}
    \item \textbf{成分の固定と自由度の導入}: $2$ 冪の法である $\mathbb{Z}_{256}$ 成分を $\mathbf{x}_0^{(256)}$ に固定し、$\mathbb{Z}_3$ 成分に摂動 $\mathbf{\delta}^{(3)}$ を加える。
    \begin{equation}
        \mathbf{x} \equiv (\mathbf{x}_0^{(3)} + \mathbf{\delta}^{(3)}, \mathbf{x}_0^{(256)})
    \end{equation}
    \item \textbf{部分空間のサンプリング}: $\mathbf{\delta}^{(3)} \in \mathbb{Z}_3^{d_a}$ を体系的に変化させる。$d_a=12$ の場合、$\mathbb{Z}_3$ 側の全空間は $3^{12} = 531,441$ 通りであり、計算機による数秒の処理で全走査または広範なサンプリングが可能である。
    \item \textbf{CRT合成（リフティング）}: 条件を満たす成分ペアを以下の写像 $\phi$ により復元する。
    \begin{equation}
        \mathbf{x} = \phi(\mathbf{x}^{(3)}, \mathbf{x}^{(256)}) = (256 \cdot y_1 \cdot \mathbf{x}^{(3)} + 3 \cdot y_2 \cdot \mathbf{x}^{(256)}) \pmod{768}
    \end{equation}
    ここで $y_1 = 256^{-1} \pmod 3$、$y_2 = 3^{-1} \pmod{256}$ である。
\end{enumerate}

\section{今後の方針}
\begin{itemize}
  \item 複数の」特殊解について、長さ8以上のサイクルの個数を比較し、性能の違いを評価する。
  \item $\aU$をランダムに生成し、それに対して整数計画法で特殊解を求め、用性能を示すものを探す。また、UTCBC以外の長さ8のサイクルを含まないものが見つけられたらうれしい。
  \item ガースが10になるような符号の構成法を考える。
\end{itemize}


\bibliography{refs}
\bibliographystyle{unsrt} %参考文献出力スタイル
\end{document}